\chapter{Implementation}
\label{chap:implementation}

Táto kapitola detailne popisuje technickú realizáciu navrhnutého systému MicroTender. Zameriava sa na vývoj smart kontraktov, integráciu decentralizovaného úložiska a tvorbu klientskej aplikácie. Hlavný dôraz je kladený na štruktúru a správanie smart kontraktu, pretože ten predstavuje nemennú logiku celého riešenia.

\section{Development environment and technologies}

Pre vývoj decentralizovanej aplikácie (dApp) sme zvolili moderný technologický stack, ktorý zabezpečuje bezpečnosť, efektivitu a dobrú používateľskú skúsenosť.

\begin{itemize}
  \item \textbf{Solidity (verzia 0.8.19):} primárny jazyk pre písanie inteligentných
  zmlúv~\cite{solidityDocs}. Verzia 0.8.x zavádza predvolenú kontrolu pretečenia
  čísel (\emph{overflow} a \emph{underflow}) pri aritmetických operáciách mimo blokov
  \texttt{unchecked}, čo zvyšuje bezpečnosť implementácie.

  \item \textbf{Hardhat:} vývojové prostredie pre Ethereum, ktoré sme využili na
  kompiláciu, automatizované testovanie a nasadzovanie zmlúv na testovaciu
  sieť~\cite{hardhatDocs}.

  \item \textbf{React (aktuálne verzia 19.x):} knižnica pre komponentové používateľské
  rozhranie aplikácie~\cite{reactDocs}. Umožňuje reaktívne aktualizácie stavu na základe
  udalostí z blockchainu (napr.\ po potvrdení transakcie) a~odpočúvania eventov kontraktu.

  \item \textbf{Ethers.js:} knižnica na komunikáciu medzi frontendom a blockchainom
  (RPC, čítanie stavu kontraktu, podpisovanie transakcií)~\cite{ethersDocs}.
  V~prehliadači sa načítava UMD zostavenie \textbf{verzie 5.7.2} z CDN
  (\texttt{Web3Provider}, \texttt{JsonRpcProvider}); v~\texttt{package.json} môže byť
  uvedená aj novšia major verzia ako závislosť pre budúce zlúčenie, no aktuálna
  implementácia \texttt{App.jsx} volá API vo formáte Ethers~v5.

  \item \textbf{Pinata a IPFS:} služba na \emph{pinning} súborov v sieti IPFS
  (\emph{InterPlanetary File System}); dokumentácia tendra sa nahráva cez REST API
  Pinata a v inteligentnej zmluve sa ukladá len identifikátor obsahu
  (CID)~\cite{pinataDocs}. Dostupnosť súborov pre používateľov závisí od pinningu,
  konfigurácie API a verejného alebo vlastného IPFS gateway; v implementácii sa na
  zobrazenie používa URL odvodené od CID (nie priamy zápis binárnych súborov do
  blockchainu).

  \item \textbf{Tailwind CSS a PostCSS:} utility-first prístup k štýlovaniu komponentov
  a spracovanie štýlov v build reťazci frontendu.

  \item \textbf{Create React App (\texttt{react-scripts}):} nástroj na lokálny vývoj
  a produkčné zostavenie frontendu bez vlastnej konfigurácie zostavovača od nuly.

  \item \textbf{Testovacia sieť Polygon Amoy:} EVM-kompatibilný reťazec, na ktorom je
  možné nasadiť zmluvu a overiť dApp s reálnymi transakciami pri nižších nákladoch
  ako na hlavnej sieti Ethereum.

  \item \textbf{Web3 peňaženka (MetaMask alebo kompatibilná):} rozšírenie prehliadača,
  ktoré používateľovi umožňuje pripojiť sa k sieti, spravovať účty a podpisovať
  transakcie volané z aplikácie cez Ethers.js.
\end{itemize}

\section{Smart contract implementation}

Jadro logiky celého systému je zapuzdrené v jednom smart kontrakte s názvom \texttt{MicroTender.sol}, napísanom v jazyku Solidity verzie 0.8.19.

\subsection{Architectural decision: Monolithic vs Factory}

Pri návrhu systému sme zvažovali použitie návrhového vzoru \textit{Factory}, kde by každý nový tender bol reprezentovaný samostatným smart kontraktom nasadeným na blockchain. Po dôkladnej analýze transakčných nákladov (Gas cost analysis) sme tento prístup prehodnotili a rozhodli sme sa pre \textbf{monolitickú architektúru}.

V tomto modeli jeden hlavný kontrakt spravuje stav a dáta všetkých tendrov pomocou dátových štruktúr. Toto rozhodnutie prináša kľúčovú výhodu v podobe výraznej úspory nákladov:
\begin{itemize}
    \item \textbf{Nasadenie nového kontraktu (Factory):} rádu $10^6$ Gas jednorazovo pri nasadení ďalšieho kontraktu na tender (presná hodnota a~ekvivalent v~ fiat závisia od siete a ceny plynu),
    \item \textbf{Vytvorenie záznamu tendra v úložisku monolitu:} výrazne nižšia jednorazová spotreba pri \texttt{createTender} oproti nasadeniu novej zmluvy.
\end{itemize}
Pre systém zameraný na "mikro" tendre je táto úspora zásadná pre jeho reálnu použiteľnosť.

\subsection{Contract structure and state variables overview}

Kontrakt \texttt{MicroTender} využíva nasledujúce kľúčové stavové premenné (Storage):

\begin{itemize}
    \item \texttt{address public owner} – vlastník kontraktu (deployer).
    \item \texttt{uint256 public tenderCounter}, \texttt{uint256 public bidCounter} – počítadlá pre generovanie unikátnych ID.
    \item \texttt{mapping(uint256 => Tender) public tenders} – hlavné úložisko tendrov indexované podľa ID.
    \item \texttt{mapping(uint256 => Bid[]) public tenderBids} – pre každé \texttt{tenderId} zoznam ponúk.
    \item \texttt{mapping(address => UserRole) public userRoles} – priradenie rolí (Member, Admin) k adresám používateľov.
    \item \texttt{mapping(address => bool) public registeredVendors} – zoznam schválených dodávateľov.
    \item \texttt{mapping(uint256 => mapping(address => bool)) public hasVoted} – záznam o tom, či daný člen už hlasoval v konkrétnom tendri.
\end{itemize}

Medzi hlavné verejné funkcie patria napr.\ \texttt{createTender},
\texttt{publishTender}, \texttt{updateTenderIPFSCID}, \texttt{submitBid},
\texttt{startVoting}, \texttt{castVote}, \texttt{finalizeTender},
\texttt{fulfillTender}, \texttt{cancelTender}, \texttt{submitVendorApplication},
\texttt{approveVendorApplication}, \texttt{rejectVendorApplication}. Na správu
rolí a dodávateľov slúžia \texttt{grantRole}, \texttt{revokeRole},
\texttt{revokeVendorStatus} a čítacie funkcie (\texttt{getTender},
\texttt{getTenderBids}, \texttt{getBidCount}, \dots). Rozsah kontraktu zodpovedá
celému životnému cyklu tendra, RBAC a validáciám bez závislosti na externých
kontraktoch.

\subsection{Constructor and initialization}

Pri nasadení kontraktu sa v konštruktore nastaví \texttt{owner} na \texttt{msg.sender}, tejto adrese sa pridelí \texttt{UserRole.Admin}, nastaví sa \texttt{hasRole[msg.sender] = true} a emituje sa \texttt{RoleGranted}. Tým má prvý nasadzovateľ oprávnenie spravovať role a systém.

\subsection{Data structures}

Pre efektívne ukladanie údajov on-chain využívame štruktúry (\texttt{struct}), ktoré zoskupujú súvisiace dáta. To umožňuje načítať všetky informácie o tendri alebo ponuke v jednom volaní.

\subsubsection{Tender structure}

Záznam o jednom tendri je reprezentovaný nasledujúcou štruktúrou:

\begin{lstlisting}[language=Solidity, caption={Štruktúra Tender v kontrakte MicroTender.sol}, label={lst:tender-struct}]
struct Tender {
    uint256 id;
    address creator;
    string title;
    string description;
    uint256 maxBudget;
    string category;
    string ipfsCID;      // Hash dokumentacie na IPFS
    uint256 deadline;
    uint256 votingDeadline;
    TenderStatus status; // Enum stavov životného cyklu (viz text)
    uint256 createdAt;
}
\end{lstlisting}

Pole \texttt{ipfsCID} uchováva identifikátor obsahu (CID) dokumentácie uloženej v IPFS; na blockchaine sa tak neukladá celý súbor, čo šetrí náklady. Stav tendra je riadený výčtovým typom \texttt{TenderStatus}. Z~pohľadu používateľského rozhrania sa prezentujú najmä fázy \textbf{zber ponúk} (\texttt{Open}), \textbf{hlasovanie} (\texttt{Voting}), \textbf{ukončenie} (\texttt{Completed}), \textbf{splnenie} (\texttt{Fulfilled}) a~\textbf{zrušenie} (\texttt{Cancelled}). V~Solidity zostáva ešte \textbf{predchádzajúci interný stav} pred prvým zverejnením (technicky \texttt{Draft}); aplikácia ho nevyužíva ako samostatný krok --- po úspešnom vytvorení tendra nasleduje bezprostredne volanie \texttt{publishTender} v rámci jedného postupu „vytvoriť a uverejniť“. Všetky tendre sú v mapovaní \texttt{mapping(uint256 => Tender) public tenders}, čo umožňuje prístup k dátam s časovou zložitosťou $O(1)$ podľa ID tendra.

\subsubsection{Bid structure}

Ponuky dodávateľov sú reprezentované štruktúrou \texttt{Bid}:

\begin{lstlisting}[language=Solidity, caption={Štruktúra Bid pre ponuku dodávateľa}, label={lst:bid-struct}]
struct Bid {
    uint256 id;
    uint256 tenderId;
    address vendor;        // Adresa dodavatela
    uint256 price;         // Ponukova cena
    uint256 deliveryTime;  // Cas dodania
    string description;
    uint256 submittedAt;
}
\end{lstlisting}

Ponuky sú evidované v dynamickom poli \texttt{Bid[]} v rámci mapovania \texttt{tenderBids}. To umožňuje jednoduché iterovanie cez všetky ponuky konkrétneho tendra počas vyhodnocovania.

\subsection{Access control}

Bezpečnosť systému je zabezpečená prostredníctvom modifikátorov funkcií (\texttt{modifiers}). Implementovali sme Role-Based Access Control (RBAC) s nasledujúcimi rolami:

\begin{itemize}
    \item \textbf{onlyOwner:} Prístup má iba vlastník kontraktu (administrátor). Používa sa na správu rolí a kritické zmeny (\texttt{grantRole}, \texttt{revokeRole}).
    \item \textbf{onlyMember:} Prístup majú iba členovia študentskej rady (a Admini, keďže \texttt{Admin >= Member}). Umožňuje vytvárať tendre, hlasovať a schvaľovať dodávateľov.
    \item \textbf{Registrovaný dodávateľ:} Prístup k funkcii \texttt{submitBid} majú iba adresy, ktoré sú v mapovaní \texttt{registeredVendors} nastavené na \texttt{true}.
\end{itemize}

Modifikátory sú definované ako funkcie, ktoré pred vykonaním pôvodnej logiky skontrolujú podmienku a v prípade nesplnenia volajú \texttt{revert}. Ukážka definície modifikátora \texttt{onlyMember}:

\begin{lstlisting}[language=Solidity, caption={Modifikátor onlyMember pre prístup členov rady}, label={lst:modifier-member}]
modifier onlyMember() {
    require(hasRole[msg.sender] && userRoles[msg.sender] >= UserRole.Member, "Only members");
    _;
}
\end{lstlisting}

Symbol \texttt{\_} označuje miesto vloženia tela pôvodnej funkcie. Registrácia členov prebieha cez funkciu \texttt{grantRole}, ktorú môže volať len \texttt{onlyOwner}.

\subsection{Events}

Smart kontrakt emituje udalosti (\texttt{event}) pri dôležitých zmenách stavu. Frontend ich môže odpočúvať cez Ethers.js a podľa nich aktualizovať UI bez potreby periodického dotazovania.

\begin{lstlisting}[language=Solidity, caption={Vybraté deklarácie eventov v MicroTender.sol}, label={lst:events}]
event TenderCreated(uint256 indexed tenderId, address indexed creator, string title);
event BidSubmitted(uint256 indexed tenderId, uint256 bidId, address vendor, uint256 price);
event VoteCasted(uint256 indexed tenderId, uint256 bidId, address voter);
event TenderCompleted(uint256 indexed tenderId, uint256 winningBidId);
\end{lstlisting}

Parametre označené \texttt{indexed} umožňujú efektívne filtrovanie logov podľa týchto hodnôt (napr. získať všetky udalosti pre daný \texttt{tenderId}). Tým sa zjednodušuje integrácia s klientskou aplikáciou a zároveň sa vytvára auditovateľná stopa na blockchaine.

\subsection{Key contract functions}

\subsubsection{Creating a tender (createTender)}

Funkcia inicializuje nový záznam v štruktúre \texttt{Tender}. Prijíma parametre ako názov, popis, rozpočet, kategóriu a \texttt{ipfsCID} (identifikátor dokumentácie v IPFS). Všetky vstupy sú validované pomocou \texttt{require} (neprázdne reťazce, rozpočet $> 0$). Funkcia využíva blok \texttt{unchecked} pre inkrementáciu počítadla \texttt{tenderCounter}, čo šetrí Gas, keďže pretečenie \texttt{uint256} je v tomto kontexte nereálne. Po zápise do \texttt{tenders} sa emituje \texttt{TenderCreated}. Zjednodušená ukážka jadra funkcie:

\begin{lstlisting}[language=Solidity, caption={Jadro funkcie createTender (validácia a zápis)}, label={lst:create-tender-code}]
function createTender(
    string memory _title,
    string memory _description,
    uint256 _maxBudget,
    string memory _category,
    string memory _ipfsCID
) external onlyMember returns (uint256) {
    require(bytes(_title).length > 0, "Title required");
    require(_maxBudget > 0, "Budget must be > 0");
    
    unchecked {
        tenderCounter++;
    }
    
    tenders[tenderCounter] = Tender({
        id: tenderCounter,
        creator: msg.sender,
        title: _title,
        description: _description,
        maxBudget: _maxBudget,
        category: _category,
        ipfsCID: _ipfsCID,
        deadline: 0,
        votingDeadline: 0,
        status: TenderStatus.Draft,
        createdAt: block.timestamp
    });
    
    emit TenderCreated(tenderCounter, msg.sender, _title);
    return tenderCounter;
}
\end{lstlisting}

Funkcia \texttt{createTender} zapíše metadáta tendra a pripraví záznam; v~záujme
jednoduchého používateľského postupu aplikácia bezprostredne volá
\texttt{publishTender} (dve transakcie za sebou pri jednom úkone „vytvoriť a
uverejniť“), čím sa tender dostane do stavu \texttt{Open} s nastaveným
\texttt{deadline} a stáva sa prístupný pre ponuky dodávateľov.

\subsubsection{Publishing a tender (publishTender)}

Funkcia \texttt{publishTender(uint256 \_tenderId, uint256 \_daysUntilDeadline)}
nastaví lehotu \texttt{deadline} a aktivuje zber ponúk (stav \texttt{Open}). Ak by
bol potrebný dodatočný zápis iného CID dokumentácie pred prvým zverejnením, kontrakt
to umožňuje cez \texttt{updateTenderIPFSCID} len v~predchádzajúcom (predzverejnenom)
stave záznamu; v~bežnom workflow aplikácie tento medzikrok používateľ nevidí.

\subsubsection{Submitting a bid (submitBid)}

Dodávateľ volá \texttt{submitBid(tenderId, price, deliveryTime, description)}. Kontrakt overí, či volajúci je v \texttt{registeredVendors}, či tender existuje a je v stave \texttt{Open}, či neprekročil \texttt{deadline} a či \texttt{price} nepresahuje \texttt{maxBudget} tendra. Nová štruktúra \texttt{Bid} sa pridá do poľa \texttt{tenderBids[tenderId]} a emituje sa \texttt{BidSubmitted}.

\subsubsection{Voting (castVote)}

Hlasovanie je povolené len v stave \texttt{Voting}. Funkcia \texttt{castVote(tenderId, bidId)} kontroluje, či volajúci je člen rady (\texttt{onlyMember}), či pre daný tender ešte nehlasoval (mapovanie \texttt{hasVoted[tenderId][msg.sender]}), či \texttt{bidId} patrí k \texttt{tenderId}. Následne sa zvýši počítadlo hlasov pre danú ponuku a označí sa, že člen už hlasoval.

\subsubsection{Starting voting (startVoting)}

Tvorca tendra môže spustiť hlasovanie len v stave \texttt{Open}. Kontrakt vyžaduje
aspoň \texttt{MIN\_BIDS\_FOR\_VOTING} (v implementácii tri) podané ponuky a dĺžku
hlasovania v rozmedzí konštánt \texttt{MIN\_VOTING\_DAYS} až \texttt{MAX\_VOTING\_DAYS}
(v implementácii 3--14 dní). Lehota na podávanie ponúk pri \texttt{publishTender}
je obmedzená rovnakým rozsahom dní (3--14).

\subsubsection{Finalizing a tender (finalizeTender)}

Funkcia \texttt{finalizeTender(tenderId)} môže byť volaná po uplynutí \texttt{votingDeadline}. Algoritmus prechádza všetky ponuky priradené k danému tendru (\texttt{tenderBids[tenderId]}), porovnáva počet hlasov a určí ponuku s najvyšším počtom hlasov za víťaznú. Tender sa prevedie do stavu \texttt{Completed} a emituje sa \texttt{TenderCompleted}.

Pri \textbf{remíze} (rovnaký počet hlasov pre viacero ponúk) zostane ako víťazná
prvá ponuka v poradí iterácie, ktorá dosiahla toto maximum --- v praxi ide o
deterministický tie-break podľa poradia v poli \texttt{tenderBids} (zvyčajne skôr
podaná ponuka). Samostatné druhé kolo hlasovania kontrakt nevykonáva.

\subsection{Gas cost optimization}

Okrem použitia \texttt{unchecked} tam, kde je to bezpečné, sme zvolili úsporný dátový model: na reťazci sú len identifikátory a metadáta, zatiaľ čo plná dokumentácia je v IPFS. Typ \texttt{uint256} je na EVM efektívny; zoskupenie dát do \texttt{struct} a prístup cez \texttt{mapping} umožňuje konštantný čas prístupu a predvídateľnú spotrebu Gas. Monolitický kontrakt namiesto factory pattern znižuje jednorazové náklady pri vytvorení tendra o rádovo 90\%.

\section{Client application implementation (Frontend)}

Klientská aplikácia slúži ako jediné používateľské rozhranie k smart kontraktu a k úložisku IPFS. Bola implementovaná vo frameworku React s použitím knižnice Ethers.js pre komunikáciu s blockchainom.

\subsection{Application structure and entry point (App.jsx)}

Projekt React aplikácie je organizovaný do priečinkov \texttt{src/},
\texttt{src/components/} (vrátane \texttt{src/components/screens/}),
\texttt{src/hooks/}, \texttt{src/utils/} a pod. \textbf{Routovanie URL sa
nepoužíva} (žiadny \texttt{react-router-dom}); navigácia medzi obrazovkami
(Dashboard, zoznamy tendrov, detail, vytvorenie tendra, hlasovanie, registrácia
dodávateľa a i.) prebieha cez stav \texttt{activeScreen} a podmienkové vykreslenie
komponentov v~\texttt{App.jsx}.

Vstupným bodom je súbor \texttt{App.jsx}, kde sú sústredené stav pripojenia
peňaženky, inštancia kontraktu, načítané tendre, vybraný tender, formuláre a
volania funkcií kontraktu. Spoločná štruktúra rozhrania (postranné menu,
hlavička, tlačidlo na pripojenie MetaMask) je rozdelená do komponentov
\texttt{Sidebar} a \texttt{Header}. Doplňujúca logika (napr. notifikácie z
eventov kontraktu) je v hooku \texttt{useNotifications}.

\subsection{Wallet connection and contract instance}

Pri pripojení peňaženky sa volá \texttt{provider.send("eth\_requestAccounts", [])}
na inštancii \texttt{ethers.providers.Web3Provider(window.ethereum)} (knižnica
Ethers~v5 načítaná z CDN), potom \texttt{getSigner()} a
\texttt{new ethers.Contract(adresa, ABI, signer)}. Bez pripojenej peňaženky
aplikácia vytvorí read-only inštanciu kontraktu cez
\texttt{ethers.providers.JsonRpcProvider} na verejný RPC uzol siete Polygon Amoy,
aby bolo možné tendre prezerať bez podpisovania transakcií.

\subsection{IPFS integration (Pinata)}

Nahrávanie dokumentácie tendra prebieha cez modul \texttt{src/utils/pinata.js}:
súbor (PDF, DOC, DOCX) sa odošle ako \texttt{multipart/form-data} na
\texttt{POST https://api.pinata.cloud/pinning/pinFileToIPFS} s autorizáciou JWT;
odpoveď obsahuje \texttt{IpfsHash} (CID), ktorý sa uloží do kontraktu v poli
\texttt{ipfsCID}. Samostatné volanie \texttt{pinJSONToIPFS} v kóde \textbf{nie je}
implementované; metadáta tendra (názov, popis, rozpočet) sú uložené on-chain.

Funkcia \texttt{getIPFSUrl} v praxi skladá URL cez verejný gateway
\texttt{https://ipfs.io/ipfs/<CID>} (v kóde je pripravená aj možnosť prepnúť na
gateway Pinata). Používateľ tak otvára dokumentáciu mimo reťazca, pričom integrita
obsahu je viazaná na CID zapísaný v kontrakte.


%%%%%%%%%%%%%%%%%%%%%%%%%%%%%%%%%%%%%%%%%%%%%%%%%%%%%%
\section{Bezpečnosť a transparentnosť}
\label{sec:security}

Navrhované riešenie je postavené na princípe, že dôležité pravidlá
obstarávania majú byť vynucované technicky, nie len organizačne.
Bezpečnosť a~transparentnosť preto netvoria samostatnú doplnkovú
funkciu systému, ale sú súčasťou jeho základného návrhu.

%%%%%%%%%%%%%%
\subsection{Bezpečnostné vlastnosti}
\label{subsec:security_properties}

Z~pohľadu bezpečnosti sú v~návrhu zohľadnené nasledujúce vlastnosti:

\begin{itemize}
    \item \textbf{Kontrola prístupu.} Prístup ku kritickým funkciám je
    riadený cez modifikátory \texttt{onlyOwner} a~\texttt{onlyMember}.
    Modifikátor \texttt{onlyOwner} obmedzuje správu rolí na vlastníka
    kontraktu, modifikátor \texttt{onlyMember} obmedzuje operácie
    spojené s~tendrami na členov rady. Dodávateľ môže podať ponuku
    len po schválení registrácie, čo je overované podmienkou
    \texttt{registeredVendors[msg.sender]}.

    \item \textbf{Validácia vstupov.} Každá verejná funkcia kontraktu
    obsahuje \texttt{require} podmienky, ktoré validujú vstupné údaje.
    Overuje sa platnosť identifikátorov cez modifikátory
    \texttt{validTender} a~\texttt{validApplication}, správnosť stavov,
    dodržanie lehôt, nenulovosť adries a~maximálna dĺžka reťazcov
    (\texttt{MAX\_STRING\_LENGTH}).

    \item \textbf{Ochrana proti pretečeniu.} Solidity vo verzii 0.8.19
    poskytuje vstavanú kontrolu pretečenia a~podtečenia celočíselných
    typov. V~kontrakte sú zároveň využité bloky \texttt{unchecked} na
    miestach, kde je pretečenie logicky vylúčené (inkrementácia
    počítadiel), čo znižuje spotrebu gas bez vplyvu na bezpečnosť.

    \item \textbf{Prevencia duplicitného hlasovania.} Mapovanie
    \texttt{hasVoted[\_tenderId][msg.sender]} zabraňuje tomu, aby
    jeden člen hlasoval v~tom istom tendri viackrát.

    \item \textbf{Prevencia duplicitnej registrácie.} Kontrakt overuje,
    že žiadateľ ešte nie je registrovaným dodávateľom a~nemá
    rozpracovanú žiadosť, čím sa zabráni opakovanému podávaniu
    žiadostí.

    \item \textbf{Obmedzenie na tvorcu tendra.} Funkcie
    \texttt{publishTender()}, \texttt{cancelTender()},
    \texttt{startVoting()}, \texttt{finalizeTender()} a~\texttt{fulfillTender()}
    sú obmedzené na adresu, ktorá tender vytvorila, prostredníctvom
    podmienky \texttt{tender.creator == msg.sender}.
\end{itemize}

Návrh zároveň zámerne neobsahuje zložitejšie bezpečnostné mechanizmy,
ako sú anonymné hlasovanie, on-chain escrow, multisig správa alebo
automatizované platby. Ich neprítomnosť nie je chybou návrhu, ale vedomým
zúžením rozsahu na jadrové vlastnosti systému, ktoré je možné v~rámci
bakalárskej práce implementovať a~otestovať spoľahlivo. Tieto rozšírenia
sú diskutované v~záverečnej kapitole ako možný smer budúcej práce.

%%%%%%%%%%%%%%
\subsection{Transparentnosť a auditovateľnosť}
\label{subsec:transparency}

Transparentnosť systému je zabezpečená tým, že kľúčové operácie sú
zapisované do blockchainu a~zanechávajú auditovateľnú stopu. Kontrakt
emituje nasledujúce udalosti (events):
\begin{itemize}
    \item \texttt{TenderCreated} -- pri vytvorení tendra,
    \item \texttt{TenderPublished} -- pri publikovaní tendra,
    \item \texttt{TenderCancelled} -- pri zrušení tendra,
    \item \texttt{BidSubmitted} -- pri podaní ponuky,
    \item \texttt{VoteCasted} -- pri odovzdaní hlasu,
    \item \texttt{TenderCompleted} -- pri finalizácii tendra,
    \item \texttt{VendorApplicationSubmitted} -- pri podaní žiadosti
    dodávateľa,
    \item \texttt{VendorApplicationApproved} -- pri schválení žiadosti,
    \item \texttt{VendorApplicationRejected} -- pri zamietnutí žiadosti,
    \item \texttt{VendorRegistered} -- pri registrácii dodávateľa (emitované pri schválení),
    \item \texttt{RoleGranted} a~\texttt{RoleRevoked} -- pri zmene rolí,
    \item \texttt{IPFSCIDUpdated} -- pri aktualizácii dokumentácie tendra.
\end{itemize}

Tieto udalosti sú uložené v~transakčných logoch blockchainu a~sú
dostupné prostredníctvom blockchain explorerov, napríklad Polygonscan
Amoy. Ľubovoľný používateľ tak môže spätne overiť, kto tender vytvoril,
kedy bol publikovaný, aké ponuky boli podané, kto hlasoval a~ktorá
ponuka zvíťazila.

Dôležitú úlohu v~transparentnosti zohráva aj prepojenie blockchainu
a~IPFS. Dokumentácia tendra nie je len pripojená ako externá príloha,
ale je reprezentovaná identifikátorom CID, ktorý umožňuje overiť, či
sa obsah dokumentu od času uloženia nezmenil. Tým sa posilňuje
dôveryhodnosť celého obstarávacieho procesu -- zmena dokumentu by viedla
k~odlišnému CID, čo by bolo okamžite zistiteľné porovnaním s~hodnotou
uloženou na blockchaine.

Výsledkom návrhu je riešenie, ktoré nie je zamerané na úplnú automatizáciu
všetkých ekonomických a~administratívnych aspektov verejného obstarávania,
ale na transparentnú a~technicky konzistentnú podporu mikro-tendrov
v~prostredí študentskej rady. Všetky dôležité udalosti sú zaznamenané
nemenným spôsobom, pravidlá sú vynucované kódom kontraktu a~dokumentácia
je prepojená s~blockchainom prostredníctvom kryptograficky overiteľných
identifikátorov.
